\section{Empirical overview: 81 conflicts, 1899--2026}\label{sec:empirical}

The contradiction-radius diagnostic requires demographic micro-data and manpower estimates, and cannot be applied sharply to every historical conflict. The 81-conflict dataset is used here only to map where civilian-share uncertainty is large, where indirect mortality dominates, and which conflicts would benefit most from identified-deaths samples; the supplementary material provides a single master table detailing the bounds, binding relations, and confidence grades for all 81 conflicts, along with a complete list of primary sources used for each.

\begin{figure}[!tp]
\centering
\includegraphics[width=0.95\textwidth]{fig_range_frame_civshare}
\caption{\textbf{Civilian share vs scale, 1899--2026.} Markers are conflicts; area increases with midpoint total deaths (min--max scaled, with a floor that keeps the smallest conflicts visible). Where sources state no ceiling for one casualty class, the class is bounded by the curated per-conflict totals or by the grand-total residual rather than treated as zero (Supplementary Material, \emph{How to read the tables}). Descriptively, the conflicts separate into two groups: combat-dominated interstate wars (civilian share $\lesssim 30\%$) and civilian-targeting events, famines, and atrocities (civilian share $\gtrsim 90\%$)---a mix of structurally different event types rather than a tested statistical bimodality.}
\label{fig:range-frame}
\end{figure}

\begin{figure}[!tp]
\centering
\includegraphics[width=0.92\textwidth]{fig_uncertainty_ladder}
\caption{\textbf{Civilian-share uncertainty ladder.} Selected high-leverage conflicts. Grey segments are the identified civilian-share interval; blue dots are the midpoints. Wide intervals mark conflicts where the next technical advance is a better identified-deaths sample, not a better plot.}
\label{fig:uncertainty-ladder}
\end{figure}

\paragraph{Generalisation.} The diagnostic transfers to any conflict with (i) a population demographic decomposition, (ii) an independent demographic survey of the dead, (iii) an independent manpower estimate at $t_0$, and (iv) a reported aggregate death total. Such inputs exist for Tigray 2020--2022 (Ethiopian Statistical Service plus WHO surveys), Syria 2011-- (Violations Documentation Center plus Syrian Network for Human Rights), the Mexican drug war (INEGI, the national statistics institute, plus CSIS/IISS cartel estimates), the Iraq War 2003--2011 \citep{burnham2006}, and Yemen 2014-- (Yemen Data Project plus the Armed Conflict Location and Event Data project, \citealp{acled}). Where these data are missing, the supplement marks the gap rather than inventing a contradiction radius.
